\section{讲解笔记：【第10次作业笔记】 力法(1)}
\begin{analysisbox}
本节按题目、答案和讲解笔记顺序整理。对原图中的结构几何关系，正文保留 可辨识文字，并在图形页中保留清理后的题图，便于核对杆件、支座、荷载和尺寸。
\end{analysisbox}
\subsection*{第 1 页转写}
1.基本概念

5-1判断题。力法方程，实际上是力的平衡方程。（X）（福州大学2012）

5-2判断题。图5-125（a）所示架结构可选用图5-125（b）所示的体系作为力

法基本体系。（V）（西南交通大学2008）

5-3

图5-126（b）所示结构可作为图5-126（a）所示结构的基本体系。（$\times$）

（烟台大学2013）

Fp

(a)

(a)

(b)

(b)

图5-125

图5-126

mmds

5-4对比图5-127（a)、（b）两个刚架的关系是（B

）。（中南大学2010）

A.内力相同，变形也相同

B.内力相同，变形不同

内力不同，变形相同

D内力不同，变形也不同

5-5图5-128所示结构的超静定次数为（/）。（中南大学2009)

A. 7

B. 8

C. 9

D. 10

2EI

4EI

EI

EI

2EI

2E1

图5.-127

图5-128

5-6图5-129（b）为图5-129（a）所示结构的力法基本体系，各杆EI=常数，k

为弹簧刚度，则其力法方程的系数、自由项和右端项分别为：81=

。（浙江大学2009)

8xit 6.p= 6) =-$\times$1/k

\subsection*{第 2 页转写}
t-6

Fp

Fp

2a

2a

(a)

(b)

图5-129

5-7判断题。图5-130所示结构中，梁AB的截面EI为常数，各链杆的EA相同，

当EI增大时，则梁截面D弯矩代数值Mp增大。（X）（中国矿业大学2011）

5-8图5-131所示两结构中，荷载Fp作用于跨中A、B处，n1、n2均为比例常数，

则当ni>n2时，跨中弯矩（

（）。（浙江大学2010)

Fp

FP

图5-130

B. MA>MB

C. MA<MB

A. MA=MB

D.不确定

\subsection*{第 3 页转写}
57

E

o

相当于EA

0

tai

EI to

相当于EA so

 0

素

 FP

LAD



t

i

鯋



5-8

i

m

篪

藏

籱

孯

f 10

t i

IOB in

Iin

j i

m二0

idf

M 0

5

-11

tiny 幽幽点

i

n



T

T

mp

Tim 分

\subsection*{第 4 页转写}
Fp

FP

图5-130

A. MA=MB

B.MA>MB

C. MA<MB

D.不确定

5-9结构上没有荷载就没有内力，这个结论在什么情况下适用？在什么情况下不适

用？（四川大学2009）

5-10判断题。图5-132所示两铰拱，已知支座B处下沉了$\triangle$B=0.01f，此时的水平

反力为零，EI=常数。（V）（天津大学2011)

5 - 11

图5-133所示结构，为了减小基础所受弯矩，可（C）。（西南交通大学

2012)

A.减小I

B.增大I2

C.增大I

D.Ii、I2同时减小n倍

EI

nEI

'n2E!

n2EI

12

12

图5-131

图5-132

图5-133

2.荷载作用下的计算

5-12用力法计算并作图5-134所示结构的M图。EI=常数。EA=6EI/?。（北京

工业大学2011)

5-13用力法计算图5-135所示结构，并作其弯矩图，设各杆EI相同。（湖南大学

2012)

5 - 14

用力法作图5-136所示结构M图。（华南理工大学2012)

5 - 15

用力法计算图5-137所示结构，并作M图。各杆EI=常数。（浙江大学2010、

烟台大学2013)

5 - 16

用力法求图5－138所示桁架支座B的反力。各杆EA相同。（山东建筑大学

2008)

5 - 17

图5-139（a）所示结构，力法计算时采用如图5-139（b）所示的基本结构：

作其M图。（中国矿业大学2010）

\subsection*{第 5 页转写}
2kN/m

3kN

10kN

q=4kN/m

EI

21

EI

1/4

6m

图5-134

图5-135

图5-136

10kN

Mo

EI

EI

10kN/m

3m

3m

(8)

(b)

图5-137

图5-138

图5-139

\subsection*{第 6 页转写}
5 2

S以40p

一訾

 二一箭



出Ti

ii

un

l

n

\_i

sn 钳 xlx t lil





cy

一  if\_

 仁i

M ynpn

5 13

8Fg\_4 4 2 10 6

i F以二

sit

一为

df\_elit.ir

2

A lj

前二从5

砉

D

如

up

0忙

5iii

811二点fixMxlxjiix535xj



2 1年2 542 15

赢

in

癕i

譽

们

in

\subsection*{第 7 页转写}
让

近

2

2乬

4  ei



1热

证2

8必

tapo

cmnnhl.lt11111111411

6

811 充 x 引 訆以6 二是2

0 针 x69 6

 zxix676

6x6

Mi

Top

一管

一15

Snxit01

x二千

no

in

i

么1

i.in hull

of

l

Mp

in

811二  xaxlxlx x4二器

op 舏ixaxnoyi ixaxno.it

一器

灶To

\subsection*{第 8 页转写}
516

gnxtopt.at

ii

o hip

o

uo

y

52

1



感

igl

砻东营

811 器x4

5  签

三  管

作檾

一听

\_\_

321臖



Snx 8以2

0忙0

\_\_\_o854

加h

针 了 X X A

Suit822 4 0770



6x xjx2

5

了

i

籇

init

iojii.ie

6Fy 10 6

1070

Fi

ftp.O

訕ixlzxitn.it

化洲

二

210

ii



Sir 等2

op

到ixzxix95xixy

izxixlx95xlitixjaii.fi

了 95成功

訫i6 45x



詍ixbx854hi

乱ix6

95X式所  x6

45 划二答

管x

一 2



o

个

一

一器

 器

\subsection*{第 9 页转写}
研究生入学考试辅导丛书结构力学（第二版）

2kN/m

3kN

10kN

=4kN/m

EI

EI

L1/4

6m

图5-134

图5-135

图5-136

10kN

Mo

EI

EI

3m

3m

[1m)

(6)

图5-137

图5-138

图5-139

3.对称性的利用

(e)

5-18图5-140所示结构中各杆EI=常数，在给定荷载作用下，H=

Ma=\_O。（浙江大学2009）

图5-141所示桁架中杆轴力FNl=

5 -19

（西南交通大学

2010)

5-20已知图5-142所示对称结构中AD杆A端的弯矩M

ql2/36（左侧受拉），

轴力FNAD= 5ql/12，则Mc=

。（下侧受拉为正）（浙江大学2010）

731

Fp/2

MBA +

10kN

al'.MBL

MBA:

HA

2m

2m

图5-141

图5-142

图5-140

230

\subsection*{第 10 页转写}
5-21判断题。图5-143所示对称结构EI=常数，中点截面C及AB杆内力应满足

M=0，Fa$\div$0，Fx=0,FNAB=0。（V）（西南交通大学2010)

5-22用力法求图5-144所示结构中CD杆的轴力，并作受弯杆件的弯矩图。设受弯

各杆EI=常数，而CD杆的EA=80。（北京交通大学2011）

5-23用力法计算图5-145所示结构，并绘出M图。已知EA=EI/16，EI=常数。

（合肥工业大学2008）

o=d+xg

g=16kN/m

EI

EA

EI

EA

2m

2m

图5-143

图5-144

图5-145

\subsection*{第 11 页转写}


ihiii

m.ae

my

ionising

齏

q

NO

it

min

4个

x

it

81x it  以2

0忙0

guy that071

0

f

mp.fr

5以

籤飈

ii

it

B

Tg

mp

op

二 Hxzx344安 管

点

 4

嶰jt.t.EE

舏  了以

这 器

ii

i

瘮

点

m

m

s

i4 咋荒

 品

811二针式44 4  t  x4 二

sun 壵fix

以以 x2

装签

二器

\subsection*{第 12 页转写}
5-24

图5-146所示结构，EI=常数，用力法作其M图（利用对称性）。（河海大学

2011)

5-25

用力法作图5-147所示结构M图，各杆EI相同。（河海大学2009）

5-26

用力法计算图5-148所示结构并作弯矩图。（西安建筑科技大学2010）

10kN/m

EI=常数

2m

28kN

28kk

EI

EI

2EI

EI

M=0

3m

4m

4m

3m

图5-146

图5-147

图5-148

5-27用力法计算图5-149所示结构，并作M图。各杆EI=常数，EA=EI/ 。（广

州大学2007)

5-28利用结构的对称性作图5-150所示结构的弯矩图。已知EI=常数，链杆EA为

无穷大。（中国矿业大学2009）

5-29

用力法作图5-151所示结构的弯矩图。EI=常数。（哈尔滨工业大学2013)

5-30

用力法分析图5-152所示刚架，绘M图。设各受弯杆EI值相同。（清华大学

2010)

21/11= va

20kN

20kN

6m

6m

图5-149

图5-150

20kN.

EA=00

3m

图5-151

图5-152

\subsection*{第 13 页转写}
25

SuXi 112  0忙0

op

it

suxittwx.tn

0 针加21 x8

2 j

8若

iiiii.nu

嵛

AT

up

AT

m

ATT4m

和 必以14 4 诈

了

千二0

s

吉ixuixz.nl  訆i 收

 譽

15MitMP

二

如

一 fix255 4 2 二

一器

18

18

18

1

出

1

64kn

h

iii

mlkninymh.li

18

8

1.5

彪

靠

靠

15

42

M

up

Snxt0p 0

F

 武泬  了咖啡等

 排M

01F  式3 4415

1器

加12KN

\subsection*{第 14 页转写}
5 27

811xitbnxz.top 0

g\_s

 叫

oil

就

I



可

吃少

ii

4

啊

咋乱批织补器

如 flxi  謍

器

加二是

in

一  吐



OY

二0

 0.07899l

mi

X2二

0.1189lxnitxuitupmf.x.to

to


To

40

in

T

h

fi

T

40 sins

in

youu

g

二专攻

  的

了

whyt.hn

0忙

一针

440 6

志 以140 120 6

2

6m

 1 if

xinitmpm

\subsection*{第 15 页转写}
5 29

i

12

节

i

i

i

i

is

Mt01Fo

贰if

ii

amp

in

Sn二 fax为2

 fax  k 管





 xii

礜

如

xnnt up M

5 30

6

It



i5 点

6

16

Inii

an

hi

Mlk

up

in

811二  x1.5x1.5x1.5xit2 141.5x1.5





一 ixl.hu xit ixzxuxl.tl



如4KN

\subsection*{第 16 页转写}
4.弹性支承超静定结构的计算

5-31图5-153（a）所示结构，取图5-153（b）为力法基本体系，其力法典型方程

为（

）。k为弹簧刚度。（中南大学2012)

Fp

Fp

20kN

EA=00

3m

(a)

(b)

图5-151

图5-152

图5-153

5-32求图5-154所示结构在给定基本体系下的力法方程及方程中的所有系数。（青

岛理工大学2008)

5-33

试计算图5-155所示具有弹性支座的结构，并作弯矩图。抗转刚度系数ks=

3EI

（中国矿业大学2008）

EI

EI

EI-常数

k-3EI/P

图5-154

图5-155

\subsection*{第 17 页转写}
532

Sixthzxz.top

一炎

duxithntopz0

tl

iiiiiici


\_29筘

\_\_\_ 器



it

in

a或以浏 点噐

如应 以的 笇意

sin

一  it 等二点

魆

0 怪着

一器

1fnnnnnn.FI

t

n

喷

811二  ixlxkjx2 訆或以州吵

it

ai 球

up

in

咋 xil 2xhxltlxz.nl

一亩tlxmxlx引一器

加M

x T T Mp M

\subsection*{第 18 页转写}
5.支座发生位移和温度变化时超静定结构的计算

5-34图5-156中等EI值单跨梁两侧温度变化时，M图位于（

（）。（青岛理工大

学2008）

B.高温侧

C. 不同侧

A.低温侧

D. 零

[0X+012X2+A1c

判断题。取图5－157所示结构基本体系，则力法方程为

5 - 35

02X+022X2+2c=2

Ibie+ Fe. c0

Ix bict 0

（福州大学2011)

0.4-0

OT

5 - 36

图5-158所示刚架除承受荷载外，支座A下沉了$\triangle$=2cm，且顺时针转动了0=

0.01rad。已知E=2.1$\times$10*kN/cm ，I=2000cm，用力法求此刚架的弯矩图。（华南理工

大学2009)

50kN

+2r

90

(b)

4m

(a)

图5-156

图5-157

图5-158

\subsection*{第 19 页转写}
5 - 36

图5-158所示刚架除承受荷载外，支座A下沉了$\triangle$=2cm，且顺时针转动了Q=

0.01rad。已知E=2.1X10*kN/cm ，I=2000cm，用力法求此刚架的弯矩图。（华南理工

大学2009)

50kN

+2r

(b)

4m

(a)

图5-156

图5-157

图5-158

\&nxi+ bp+ bicc 0

50kN

b3.04

50kN

4wkn.

bp=-(xnx lww.$\times$ 4m)=

Ix blct

fe: L=0, Ix On +IxD +umx b=0

IL= -b - 4mx B= -$\times$ 1b?\_ 4$\times$0.0] = - 0.06 m

9 =IX

\subsection*{第 20 页转写}
6.超静定结构的位移计算

5-37图5-159所示梁（跨长1）支座位移产生的跨中点竖向位移为（

（哈尔

滨工业大学2011)

A. lp/16

B. lp/8

C. lp/4

D. lp/2

5-38图5-160所示为计算得到的超静定刚架的弯矩图，则C点的角位移pc=

，E点的水平位移$\triangle$HE=

（湖南大学2011）

5-39图5-161所示一端固定、另一端定向支承的等截面梁，当A端发生单位转角

9=1时，B端的竖向位移$\triangle$=

（重庆大学2010）

111.98

7.29

3.65

EI

8.33

.0

EI

2EI

EI

9=1

EI

3.65

F0.52

4.17

M图(x$\times$0.01Fpl)

图5-159

图5-160

图5-161

5-40图5-162所示结构，各杆抗弯刚度EI为常数，忽略杆件的轴向变形。（1）根

\subsection*{第 21 页转写}
研究生入学考试辅导丛书结构力学（第二版）

据结构与荷载的特点，选用最简便的方法作刚架的弯矩图；（2）求点的水平位移。（武汉

理工大学2011)

10kN

20kN-m

20kN-m

2m

2m

2m

图5-162

(()

